% 图 4-2 时序收敛过程（WNS 分阶段）
\begin{tikzpicture}
\begin{axis}[
  width=150mm, height=56mm,
  ybar, bar width=13mm,
  ymin=-1.05, ymax=0.22, xmin=0.4, xmax=4.6,
  xtick={1,2,3,4}, xticklabels={基线,策略扫掠,结构改造,板级签核},
  xticklabel style={font=\footnotesize},
  yticklabel style={font=\footnotesize},
  ylabel={WNS / ns}, ylabel style={font=\footnotesize},
  ymajorgrids, grid style={macline!25},
  axis lines*=left,
  enlarge x limits=0.12,
  nodes near coords, nodes near coords align={vertical},
  every node near coord/.append style={font=\scriptsize},
  point meta=explicit symbolic,
]
\addplot[fill=baryellow, draw=baryedge] coordinates {
  (1,-0.942) [$-0.942$]
  (2,-0.118) [$-0.118$]
  (3,0.007)  [$+0.007$]
  (4,0.087)  [$+0.087$]
};
\draw[black!55, thick, dashed] (axis cs:0.4,0) -- (axis cs:4.6,0);
\end{axis}
\end{tikzpicture}

\vspace{1mm}
{\scriptsize
\begin{tabular}{@{}p{0.16\textwidth}p{0.36\textwidth}p{0.42\textwidth}@{}}
\toprule
阶段 & 动作 & 关键读数 \\
\midrule
基线 & HEAD 默认实现策略 & TNS $-433.585$ ns，失败端点 1\,579 \\
策略扫掠 & 33 组内置策略逐一运行（并行），最优 \texttt{Performance\_ExploreWithRemap} & TNS $-1.008$ ns，失败端点 22；合计回收 $0.824$ ns（87\%） \\
结构改造 & 三处归因改造：\texttt{ascon\_ise} 复位局部化（93 端点）、\texttt{mult.v} 输入输出寄存与握手（48 端点）、\texttt{dcache} 增 \texttt{MISS\_CAPTURE}（11 端点） & 双策略 STA：$+0.001$ / $+0.007$ ns，失败端点 0；资源代价 $+118\sim+190$ LUT、$+43\sim+91$ FF \\
板级签核 & 三镜像最终实现 & TNS 0、失败端点 0（共 55\,572）、WHS $+0.021$ ns、DRC 0 Error \\
\bottomrule
\end{tabular}}
